\textbf{Program 1: Student records from definition to update.}

\textbf{File:} \texttt{student\_records.cpp}.

\textbf{Chapter connection:} C++ practical: structures; Putting records into an array.

Keep the Student fields from the chapter. Separate lookup from mutation: findById returns an index; updateStudent checks it before forming a pointer. main builds five records, searches, calculates an average, and verifies that only the selected record changes.

\begin{adslabfigure}[How this program runs]
\begin{adslabgraphic}
\begin{tikzpicture}[font=\small,execute at begin node={\hyphenpenalty=10000\exhyphenpenalty=10000}]
\node[utpstep,text width=2.50cm,minimum height=1.15cm] (s0) at (0.0,0) {\textbf{1}\\Student array};
\node[utpstep,text width=2.50cm,minimum height=1.15cm] (s1) at (2.95,0) {\textbf{2}\\findById / display};
\node[utpstep,text width=2.50cm,minimum height=1.15cm] (s2) at (5.9,0) {\textbf{3}\\averageAge};
\node[utpstep,text width=2.50cm,minimum height=1.15cm] (s3) at (8.85,0) {\textbf{4}\\updateStudent / checks};
\draw[utparrow] (s0.east)--(s1.west);
\draw[utparrow] (s1.east)--(s2.west);
\draw[utparrow] (s2.east)--(s3.west);
\end{tikzpicture}
\end{adslabgraphic}
\end{adslabfigure}


\begin{utpsteps}\setcounter{enumi}{0}\item \textbf{Headers and initial definitions.}\end{utpsteps}

\texttt{student\_records.cpp} --- lines 10--19.

\lstinputlisting[firstline=10,lastline=19]{all_experiments/Lab01/student_records.cpp}


\tcbbreak

\begin{utpsteps}\setcounter{enumi}{1}\item \textbf{Implement display, findById, updateStudent and helpers.}\end{utpsteps}

\texttt{student\_records.cpp} --- lines 20--54.

% Worked lab input; chapter fonts, colours and TikZ styles retained.
\begin{utpcompactlabfigure}[One address, two field changes]
\begin{labgraphic}
\begin{tikzpicture}[font=\small]

\node[utpstep,align=left,text width=3.7cm] (a) at (0,0) {\textbf{students[3]}\\ID: 1004 \quad Devi\\age: 22 \quad course: IS};
\node[utpstep,align=left,text width=3.7cm,draw=UTPGreen,fill=UTPPaleTeal] (b) at (8,0) {\textbf{students[3]}\\ID: 1004 \quad Devi\\age: 23 \quad course: CS};
\draw[utparrow] (a)--node[utplabel,above,align=center]{\texttt{selected->age = 23}\\\texttt{selected->course = "CS"}}(b);
\node[utplabel,align=center] at (4,-1.3) {\texttt{selected = \&students[3]} points to the original record\\The other four records remain unchanged};
\end{tikzpicture}
\end{labgraphic}
\end{utpcompactlabfigure}


\lstinputlisting[firstline=20,lastline=54]{all_experiments/Lab01/student_records.cpp}


\tcbbreak

\begin{utpsteps}\setcounter{enumi}{2}\item \textbf{Connect the test cases in main.}\end{utpsteps}

\texttt{student\_records.cpp} --- lines 55--88.

\textbf{Trace checkpoint for the supplied main.} Before the update, save the five records exactly as main does. ID 1003 is at index 2; ID 9999 must return -1. Updating ID 1004 changes only the fourth record to age 23 and course CS. Trace the guard for a missing ID before forming a pointer: an index of -1 must never be used as an array offset. Compare every record with before, not only the updated one.

\lstinputlisting[firstline=55,lastline=88]{all_experiments/Lab01/student_records.cpp}


\tcbbreak

\begin{utpsteps}\setcounter{enumi}{3}\item \textbf{Connect the test cases in main.}\end{utpsteps}

\texttt{student\_records.cpp} --- lines 89--113. Continue in the same file; do not add another header or opening brace.

\lstinputlisting[firstline=89,lastline=113]{all_experiments/Lab01/student_records.cpp}


\tcbbreak

\textbf{Build, run, and inspect student\_records.}

From this program\textquotesingle s folder, run these commands in PowerShell:

\begin{lstlisting}[style=utpterminal,language={}]

g++ -std=c++17 -Wall -Wextra student_records.cpp -o student_records.exe

.\student_records.exe

\end{lstlisting}

\textbf{Expected output excerpt.} Compare values, order, and pass/fail messages; spacing need not match.

\begin{lstlisting}[style=utpterminal,language={}]

== All records (ID | name | age | course) ==
1001 | Amina | 20 | CS
1002 | Ben | 21 | IT
1003 | Chen | 20 | CS
1004 | Devi | 22 | IS
1005 | Ehsan | 21 | IT

== Search by ID ==
Search 1003: found at index 2 -> Chen
Search 9999: not found

== Average age (a 1D-array statistic) ==
Average age: 20.8

== Update record 1004 through a pointer ==

\end{lstlisting}

The complete captured run is in \texttt{labs/outputs/student\_records.txt}.

\begin{tryit}

Change the target ID to 9999. Predict the failed search and update: both must leave every record unchanged. Then update ID 1004 and calculate the new average, 21.0, from the stored ages.

\end{tryit}
